\section*{Capítulo \ref{ch:g10:reffunc} --- Funciones de referencia}

\begin{solution}{exo:g10:reffunc:1}
$f $: \omterm{def:g10:reffunc:affine}{pendiente} $3$, interceptar $-2$, \omterm{def:g10:functions:variations}{creciente}. \quad
$ g $: \omterm{def:g10:reffunc:affine}{pendiente} $-\frac12$, interceptar $4$, \omterm{def:g10:functions:variations}{decreciente}. \quad
$ h $: \omterm{def:g10:reffunc:affine}{pendiente} $0$, interceptar $7$, constante. \quad
$ k $: \omterm{def:g10:reffunc:affine}{pendiente} $-1$, interceptar $0$, \omterm{def:g10:functions:variations}{decreciente} (y lineal).
\end{solution}

\begin{solution}{exo:g10:reffunc:2}
A través de $(2,1)$ y $(5,10)$: \omterm{def:g10:reffunc:affine}{pendiente} $ m = \frac{10 - 1}{5 - 2} = 3$; entonces
$1 = 3 \times 2 + p $ da $ p = -5$: $ f(x) = 3x - 5$.

A través de $(-1,4)$ y $(3,-4)$: \omterm{def:g10:reffunc:affine}{pendiente} $ m = \frac{-4 - 4}{3 - (-1)} =
\frac{-8}{4} = -2$; entonces $4 = -2 \times (-1) + p $ da $ p = 2$:
$ f(x) = -2x + 2$.
\end{solution}

\begin{solution}{exo:g10:reffunc:3}
$(3.1)^2 < (3.2)^2$: el cuadrado aumenta en positivos.

$(-2.7)^2 < (-2.8)^2$: el cuadrado disminuye en negativos y
$-2.8 < -2.7$.

$\frac{1}{5.1} > \frac{1}{5.2}$: lo inverso disminuye en los positivos.

$\sqrt{17} > \sqrt{15}$: la raíz cuadrada aumenta.
\end{solution}

\begin{solution}{exo:g10:reffunc:4}
$ x^2 = 16$: dos soluciones, $ x = 4$ y $ x = -4$ (la línea horizontal
$ y = 16$ corta la parábola dos veces).

$ x^2 \leq 16$: la parábola está debajo de la línea entre los dos
\omterm{def:g10:numbers:interunion}{intersección} agujas: $ x \in \intcc{-4}{4}$.

$ x^2 > 9$: la parábola está estrictamente por encima $ y = 9$ afuera
$\intcc{-3}{3}$: $ x \in \intoo{-\infty}{-3} \cup \intoo{3}{+\infty}$.
\end{solution}

\begin{solution}{exo:g10:reffunc:5}
Para $0 < x < 1$: multiplicando $ x < 1$ repetidamente por $ x $ da
$ x^3 < x^2 < x $, y \cref{prop:g10:reffunc:compare} da $ x < \sqrt x $,
entonces
\[
x^3 < x^2 < x < \sqrt x .
\]
Consulte con $ x = 0.25$: $ x^3 = 0.015625$, $ x^2 = 0.0625$, $ x = 0.25$,
$\sqrt x = 0.5$.
\end{solution}

\begin{solution}{exo:g10:reffunc:6}
$\frac1x = 2$ tiene la solución única $ x = \frac12$ (la hipérbola se encuentra
la línea horizontal $ y = 2$ una vez, en la rama positiva).

Para $ x > 0$: $\frac1x < 2$. Multiplicando por $ x > 0$: $1 < 2x $, entonces
$ x > \frac12$: conjunto de soluciones $\intoo{\frac12}{+\infty}$.

Si $ x < 0$ está permitido: todo negativo $ x $ satisface $\frac1x < 0 < 2$,
entonces el conjunto completo de soluciones es
$\intoo{-\infty}{0} \cup \intoo{\frac12}{+\infty}$.
\end{solution}

\begin{solution}{exo:g10:reffunc:7}
(a) En $\intcc{2}{5}$, todas las entradas son positivas y al elevar al cuadrado aumenta:
$4 \leq x^2 \leq 25$.

(b) en $\intcc{-3}{-1}$, al cuadrado disminuye:
el cuadrado más grande proviene de $-3$: $1 \leq x^2 \leq 9$.

(c) en $\intcc{-2}{3}$, el \omterm{def:g10:numbers:interval}{intervalo} contiene $0$, donde la plaza
alcanza su \omterm{def:g10:functions:extrema}{mínimo} $0$; el \omterm{def:g10:functions:extrema}{máximo} es el mayor de $(-2)^2 = 4$ y
$3^2 = 9$. Entonces $0 \leq x^2 \leq 9$.
\end{solution}

\begin{solution}{exo:g10:reffunc:8}
Primera empresa: $ f(x) = 1.5x + 4$; segundo: $ g(x) = 2x $. El primero es
más barato cuando
\[
1.5x + 4 < 2x
\ \Longleftrightarrow\
4 < 0.5x
\ \Longleftrightarrow\
x > 8 .
\]
Más allá de $8$ km, la primera empresa es más barata; exactamente en $8$ km ambos cargan
$16$.
\end{solution}

\begin{solution}{exo:g10:reffunc:9}
Ambos lados son no negativos, por lo que la desigualdad es equivalente a la
entre sus cuadrados:
\[
\left(\sqrt{a+b}\right)^2 = a + b
\qquad\text{and}\qquad
\left(\sqrt a + \sqrt b\right)^2 = a + 2\sqrt a \sqrt b + b .
\]
Desde $2\sqrt a\sqrt b \geq 0$, el segundo cuadrado es al menos el primero,
lo que prueba $\sqrt{a+b} \leq \sqrt a + \sqrt b $. La igualdad se cumple exactamente
cuando $\sqrt a \sqrt b = 0$, es decir,\ cuando $ a = 0$ o $ b = 0$.
\end{solution}

\begin{solution}{exo:g10:reffunc:10}
\emph{1.} Pon el lado derecho sobre un denominador común:
\[
2 + \frac{3}{x-1} = \frac{2(x-1) + 3}{x - 1} = \frac{2x + 1}{x - 1}
= f(x).
\]

\emph{2.} En $\intoo{1}{+\infty}$, como $ x $ aumenta, $ x - 1$ aumenta
y se mantiene positivo, así que $\frac{3}{x-1}$ disminuye (inversa \omterm{def:g10:functions:function}{función}, veces
la constante positiva $3$), entonces $ f(x) = 2 + \frac{3}{x-1}$ disminuye:
$ f $ es estrictamente \omterm{def:g10:functions:variations}{decreciente} en $\intoo{1}{+\infty}$.
\end{solution}

\begin{solution}{pb:g10:reffunc:1}
\textbf{1.} $ d(30) = 3^2 = 9$ metro; $ d(50) = 25$ metro; $ d(90) =
81$ metro; $ d(130) = 169$ metro.

\textbf{2.} Duplicación $ v $ dobles $\frac v{10}$ y multiplica
su cuadrado por $4$: $ d(100) = 100$ metro $= 4 \times d(50)$. dos veces
la velocidad, cuatro veces la distancia --- el cuadrado \omterm{def:g10:functions:function}{función}'s
escalado.

\textbf{3.} $\left(\frac{v}{10}\right)^2 = 50$ da
$\frac{v}{10} = \sqrt{50}$, entonces $ v = 10\sqrt{50} \approx 71$
km/h. La raíz cuadrada respondió; la lectura es inequívoca
porque las velocidades son positivas, y en
$\intco{0}{+\infty}$ la plaza \omterm{def:g10:functions:function}{función} sube estrictamente
(\cref{prop:g10:reffunc:sqrt}: uno positivo \omterm{def:g10:functions:function}{preimagen}).

\textbf{4.} $ d(31) - d(30) = 9.61 - 9 = 0.61$ m, mientras
$ d(91) - d(90) = 82.81 - 81 = 1.81$ m: lo mismo $+1$ costos km/h
tres veces más distancia a alta velocidad. La parábola no sólo
se eleva, se \emph{se intensifica} --- su tasa de ascenso crece con $ x $.

\textbf{5.} $ D(50) = 14 + 25 = 39$ metro; $ D(130) = 36.4 + 169 =
205.4$ metro. En la ciudad el término afín (reacción) es el más grande
compartir; en la carretera el término cuadrado la aplasta --- afín
crece constantemente, los cuadrados se escapan.

\textbf{6.} $ F(0.5) = 240$, $ F(1) = 120$, $ F(3) = 40$.

\textbf{7.} lo inverso \omterm{def:g10:functions:function}{función} disminuye: cuanto más lejos
el pivote, menor será la fuerza necesaria --- con un tiempo suficientemente largo
brazo de palanca ($ d $ enorme), cualquier fuerza, por pequeña que sea, equilibra cualquier
carga: el alarde de Arquímedes vive en la cola de la hipérbola, que
enfoques $0$ sin llegar a él. Y $ d = 0$ es el
Valor prohibido: sin brazo, sin palanca --- la asíntota vertical.

\textbf{8.} $\frac{100}{4} = 25$ en $2$ metro,
$\frac{100}{9} \approx 11$ en $3$ metro, $1$ en $10$ metro. a distancia
$ d $ la luz se ha extendido sobre una esfera de área $4\pi d^2$
(Problema de la lápida de Arquímedes, en el volumen de Escuela Secundaria): la misma energía dividida por un área
creciendo como $ d^2$ --- de ahí lo inverso \emph{cuadrado}.

\textbf{9.} $60$, $30$, $20$ euros per cápita por $ n = 10, 20,
30$. Para $\frac{600}{n} \leq 25$: $ n \geq 24$ participantes. el
La hipérbola promete acciones cada vez más baratas $ n $ crece --- y
se niega a alcanzar $0$: el autobús nunca es gratis.

\textbf{10.} $ a(x) = \frac{2x + 3}{x} = 2 + \frac3x $:
\omterm{def:g10:functions:variations}{decreciente} en $\intoo{0}{+\infty}$ (inverso \omterm{def:g10:functions:function}{función} desplazado),
acercándose a la asíntota $ y = 2$. Económicamente: difundir la
fijo $3$ Los euros por más carteles hacen bajar el coste unitario.
hacia lo incompresible $2$ euros de materiales --- economías
de escala, con piso duro.

\textbf{11.} Tierra: $ a^3 = 1$, $ T^2 = 1$. Marte:
$1.52^3 \approx 3.51$ y $1.88^2 \approx 3.53$. Júpiter:
$5.20^3 \approx 140.6$ y $11.86^2 \approx 140.7$. Kepler
aviso: $ T^2 = a^3$ para cada planeta --- una ley para el todo
cielo.

\textbf{12.} $ T = \sqrt{a^3}$: para Saturno
$\sqrt{9.54^3} = \sqrt{868} \approx 29.5$ años, contra el
observado $29.4$: la ley se mantiene en una décima parte.

\textbf{13.} Asteroide: $ T = \sqrt{4^3} = \sqrt{64} = 8$ años.
Cometa: $ a^3 = 27^2 = 729 = 9^3$, entonces $ a = 9$ AU.

\textbf{14.} Cuadrado ($ T^2$), cubo ($ a^3$) y raíz cuadrada (a
extraer $ T $). Escalar: multiplicar $ a $ por $4$ multiplica $ T $ por
$4\sqrt4 = 8$ --- como lo confirma el asteroide ($ a $ cuatro veces
de la Tierra, $ T $ ocho veces).

\textbf{15.} Una tabla de números, comparada con las formas de
algunas referencias funciones, produjo una ley del universo
setenta años antes de que alguien supiera \emph{por qué} se mantuvo ---
reconociendo la \omterm{def:g10:functions:function}{función} es la mitad de la ciencia.

\textbf{16.} En $\intoo{0}{1}$: $ x^2 < x < \sqrt x $; en
$\intoo{1}{+\infty}$: $\sqrt x < x < x^2$. Cheques: en
$ x = 0.25$: $0.0625 < 0.25 < 0.5$; en $ x = 4$:
$2 < 4 < 16$.

\textbf{17.} Para $ x \geq 0$, ambos lados son $\geq 0$, entonces
$\sqrt x > x \iff x > x^2 \iff x(1 - x) > 0 \iff
x \in \intoo{0}{1}$.

\textbf{18.} En $ x = 0.5$: $ x^3 = 0.125 < x^2 = 0.25 < x =
0.5$. En $ x = 2$: $2 < 4 < 8$: el orden se invierte. cruces de
$ x^2$ y $ x^3$: $ x^3 - x^2 = x^2(x - 1) = 0$: en $ x = 0$ y
$ x = 1$.

\textbf{19.} $10t^2 = 100\sqrt t $ da $ t^2 = 10\sqrt t $;
elevando al cuadrado, $ t^4 = 100t $, entonces $ t^3 = 100$ y
$ t = \sqrt[3]{100} \approx 4.64$ años --- aproximadamente $4$ años y
$8$ meses. (Ambos planes pagan aproximadamente $215$ euros.) Moraleja:
Las raíces corren temprano, los poderes ganan cada maratón.

\textbf{20.} Frenado: la parábola \emph{empinamiento} es el
Peligro de velocidad. Palanca: la \emph{asíntota vertical} en $0$
y la cola larga son la ventaja del mecánico. Lámpara: la
$ d^2$ del \emph{esfera en expansión} establece el desvanecimiento. Kepler:
el \emph{exponente} $\frac32$ es el ritmo del sistema solar.
Ahorro A: el de la raíz \emph{aplastamiento} es del ahorrador lento
destino. Cinco \omterm{def:g10:functions:graph}{gráficos}, cinco leyes: en efecto, el alfabeto de la naturaleza.
\end{solution}
